\section*{Summary}
We built binary trees whose routers and leaves are small multilayer perceptrons, fed them features from three kinds of trunk, and studied them on CIFAR-10 and Cats vs Dogs. The main results:
\begin{enumerate}
\item \emph{Overfitting (E1).} All ten trees (depth 3, 0.3M to 1.3M parameters) reached 100\% training accuracy on the full training sets, with or without the routers made hard. Test accuracy was set by the trunk: 59\%/74\% (CIFAR-10/Cats vs Dogs) with random features, 82\%/91\% with frozen ImageNet ResNet-18 features and 85\%/92\% with a CNN trained separately on the task. Routers that read only shallow trunk features cost at most 1.3 points while halving the parameters used per image.
\item \emph{Double descent (E2).} With 2{,}000 training images and the squared loss, test error peaks where the training loss reaches zero, at $P\approx nK=20$k parameters on CIFAR-10, and falls again past it (20\% noise: 46\% best small tree, 67\% at the peak, 51\% at 308k parameters). Over a grid of 40 depth--width combinations the peak follows the line $P\approx nK$, whether the parameters come from more leaves or wider leaves. More data can hurt (sample-wise peak at $n=P/K$). Cross-entropy moves the peak to about 3k parameters; weight decay $10^{-2}$ removes it and gives the best test error of the sweep (32\%).
\item \emph{Receptive fields (E3).} Every node looks at the whole image, at every depth; nodes differ in the image properties they respond to. With ImageNet features the Cats vs Dogs routers ignore the class and sort photographs by setting (cages, blankets, grass, dark fur); with the task-trained CNN, routing becomes more semantic, saliency moves to the object, and synthesised images show object parts.
\item \emph{Minimising size (E4).} A linear classifier on the same features (9.6k / 1.9k parameters) matches the 1.2M-parameter tree; distilled from the tree it is slightly better. Hard routing runs 18\% of the tree per image; 90\% magnitude pruning costs about 1 point; on Cats vs Dogs the tree can be collapsed from 8 to 3 leaves with no loss of accuracy.
\item \emph{Mislabelled images (E5).} The tree does not quarantine them in a few leaves. When leaves are tiny, the routers bend to the wrong labels; when leaves are large, the leaves memorise them locally, and the second descent comes from test images whose nearest training image is correctly labelled.
\end{enumerate}
